\chapter{§4.5 概率初步}
\label{sec:4.5}


\prerequisite{§1.5 分数, §1.7 百分数}

\gradelevel{7-9 年级}


生活中有些事情是确定的（太阳从东方升起），有些是不确定的（明天会不会下雨）。\textbf{概率}就是用数学方法来度量不确定事件发生的可能性大小。本章从随机事件的分类出发，学习如何计算概率，以及如何用列表法和树形图系统地列举所有可能的结果。



\section*{随机事件、必然事件与不可能事件}


\begin{explore}


请判断以下事件是否一定会发生：

\begin{enumerate}
  \item 掷一枚骰子，点数小于 $7$ 。
  \item 掷一枚骰子，点数等于 $8$ 。
  \item 掷一枚骰子，点数等于 $3$ 。
\end{enumerate}


\end{explore}

\begin{understand}


按照发生的确定性，事件可以分为三类：

\textbf{必然事件}：在一定条件下\textbf{一定会发生}的事件。
\begin{itemize}
  \item 例：掷骰子，点数小于 $7$ 。（骰子只有 $1 \sim 6$ ，所以必定小于 $7$ 。）
\end{itemize}


\textbf{不可能事件}：在一定条件下\textbf{一定不会发生}的事件。
\begin{itemize}
  \item 例：掷骰子，点数等于 $8$ 。（骰子最大就是 $6$ ，不可能出现 $8$ 。）
\end{itemize}


\textbf{随机事件}：在一定条件下\textbf{可能发生也可能不发生}的事件。每次实验前无法确定结果。
\begin{itemize}
  \item 例：掷骰子，点数等于 $3$ 。（可能是 $3$ ，也可能不是。）
\end{itemize}


\textbf{约定}：
\begin{itemize}
  \item 必然事件的概率 $= 1$ 。
  \item 不可能事件的概率 $= 0$ 。
  \item 随机事件的概率在 $0$ 与 $1$ 之间，即 $0 < P < 1$ 。
  \item 一个事件越容易发生，概率越接近 $1$ ；越不容易发生，概率越接近 $0$ 。
\end{itemize}


\end{understand}

\begin{workedexamples}


\textbf{例 1.} 判断下列事件属于哪一类（必然事件、不可能事件或随机事件）：

（1）在标准大气压下，水加热到 $100°\text{C}$ 时沸腾。
（2）任意买一张电影票，座位号是偶数。
（3）抛一枚硬币，正面朝上。
（4）经过一个路口，遇到红灯。
（5）在只装有白球的袋中摸出红球。

\textbf{解：}

（1）必然事件（物理规律决定）。
（2）随机事件（可能是偶数也可能是奇数）。
（3）随机事件（可能正面也可能反面）。
（4）随机事件（可能是红灯、绿灯或黄灯）。
（5）不可能事件（袋中没有红球）。

\end{workedexamples}

\begin{keytakeaway}


\begin{itemize}
  \item 事件分三类：必然事件（一定发生）、不可能事件（一定不发生）、随机事件（可能发生）。
  \item 概率 $P$ 满足 $0 \leq P \leq 1$ 。
\end{itemize}


\end{keytakeaway}



\section*{概率的意义——频率的稳定值}


\begin{explore}


抛一枚硬币，正面朝上的概率是 $\frac{1}{2}$ 。这是什么意思？是说"抛 2 次，一定有 1 次正面"吗？

当然不是！你完全可能抛 2 次都是反面。那 $\frac{1}{2}$ 到底意味着什么？

\end{explore}

\begin{understand}


历史上，许多数学家做过大量的抛硬币实验：

\begin{center}
\begin{tabular}{llll}
\toprule
实验者 & 抛掷次数 & 正面次数 & 正面频率 \\
\midrule
德·摩根 & 2048 & 1061 & $0.5181$ \\
蒲丰 & 4040 & 2048 & $0.5069$ \\
皮尔逊 & 12000 & 6019 & $0.5016$ \\
皮尔逊 & 24000 & 12012 & $0.5005$ \\
\bottomrule
\end{tabular}
\end{center}


\textbf{规律}：随着实验次数增加，正面朝上的\textbf{频率}越来越稳定地接近 $0.5$ 。

\textbf{概率的统计定义}：在大量重复实验中，随机事件 $A$ 发生的频率会稳定在一个固定的值附近，这个值就是事件 $A$ 的\textbf{概率}，记作 $P(A)$ 。

\[
P(A) = \lim_{\text{实验次数} \to \text{很大}} \frac{\text{事件 } A \text{ 发生的次数}}{\text{实验总次数}}
\]


\textbf{重要区分}：
\begin{itemize}
  \item \textbf{频率}是实验得到的结果，每次实验可能不同。
  \item \textbf{概率}是固定不变的理论值，频率在它附近波动。
  \item 实验次数越多，频率越接近概率。
\end{itemize}


\end{understand}

\begin{keytakeaway}


\begin{itemize}
  \item 概率是大量实验中频率趋近的\textbf{稳定值}。
  \item 频率可以作为概率的\textbf{估计值}，实验次数越大，估计越准。
  \item 概率不意味着"必然按比例出现"——抛 10 次硬币不一定恰好 5 次正面。
\end{itemize}


\end{keytakeaway}



\section*{等可能事件的概率（古典概型）}


\begin{explore}


掷一枚均匀骰子，出现每个点数的机会是一样的。那么，掷出点数为偶数的概率是多少？

\end{explore}

\begin{understand}


如果一个实验满足以下两个条件：
\begin{enumerate}
  \item 所有可能的结果是\textbf{有限个}。
  \item 每个结果出现的可能性\textbf{相等}（等可能）。
\end{enumerate}


则称为\textbf{等可能事件}（也叫\textbf{古典概型}）。

\textbf{等可能事件的概率公式}：

\[
P(A) = \frac{A \text{ 包含的等可能结果数}}{\text{所有等可能结果总数}}
\]


\textbf{掷骰子，点数为偶数}：
\begin{itemize}
  \item 所有等可能结果： $\{1, 2, 3, 4, 5, 6\}$ ，共 $6$ 个。
  \item 事件"点数为偶数"包含的结果： $\{2, 4, 6\}$ ，共 $3$ 个。
\end{itemize}


\[
P(\text{偶数}) = \frac{3}{6} = \frac{1}{2}
\]


\end{understand}

\begin{workedexamples}


\textbf{例 2.} 一个不透明的袋子里装有 $4$ 个红球和 $6$ 个白球，它们除颜色外完全相同。随机摸出 $1$ 个球，求：
（1）摸到红球的概率。
（2）摸到白球的概率。

\textbf{解：}

总球数 $= 4 + 6 = 10$ 个，每个球被摸到的机会相等。

（1） $P(\text{红球}) = \frac{4}{10} = \frac{2}{5}$

（2） $P(\text{白球}) = \frac{6}{10} = \frac{3}{5}$

验证： $\frac{2}{5} + \frac{3}{5} = 1$ ✓（摸出的球要么是红球要么是白球。）

\textbf{例 3.} 一副扑克牌去掉大小王后剩 $52$ 张。随机抽取 $1$ 张，求：
（1）抽到红心的概率。
（2）抽到 $\text{A}$ （四种花色）的概率。
（3）抽到红心 $\text{A}$ 的概率。

\textbf{解：}

（1）红心有 $13$ 张。 $P(\text{红心}) = \frac{13}{52} = \frac{1}{4}$

（2） $\text{A}$ 有 $4$ 张（四种花色各一张）。 $P(\text{A}) = \frac{4}{52} = \frac{1}{13}$

（3）红心 $\text{A}$ 只有 $1$ 张。 $P(\text{红心 A}) = \frac{1}{52}$

\end{workedexamples}

\begin{keytakeaway}


\begin{itemize}
  \item 等可能事件的概率 $= \frac{\text{有利结果数}}{\text{全部结果数}}$ 。
  \item 使用此公式的前提是所有基本结果\textbf{等可能}。
  \item $P(A) + P(\bar{A}) = 1$ ，其中 $\bar{A}$ 表示事件 $A$ 不发生（即 $A$ 的对立事件）。
\end{itemize}


\end{keytakeaway}



\section*{用列表法求概率}


\begin{explore}


同时掷两枚骰子，两枚骰子点数之和等于 $7$ 的概率是多少？

两枚骰子各有 $6$ 种结果，总共有 $6 \times 6 = 36$ 种组合。但哪些组合的和等于 $7$ ？逐一列出比凭直觉更可靠。

\end{explore}

\begin{understand}


\textbf{列表法}（也叫表格法）是将两个实验的所有可能组合列成一张表格，然后数出满足条件的组合个数。

\textbf{掷两枚骰子的所有结果}：

\begin{center}
\begin{tabular}{lllllll}
\toprule
骰子 1 \ 骰子 2 & 1 & 2 & 3 & 4 & 5 & 6 \\
\midrule
\textbf{1} & (1,1) & (1,2) & (1,3) & (1,4) & (1,5) & (1,6) \\
\textbf{2} & (2,1) & (2,2) & (2,3) & (2,4) & (2,5) & (2,6) \\
\textbf{3} & (3,1) & (3,2) & (3,3) & (3,4) & (3,5) & (3,6) \\
\textbf{4} & (4,1) & (4,2) & (4,3) & (4,4) & (4,5) & (4,6) \\
\textbf{5} & (5,1) & (5,2) & (5,3) & (5,4) & (5,5) & (5,6) \\
\textbf{6} & (6,1) & (6,2) & (6,3) & (6,4) & (6,5) & (6,6) \\
\bottomrule
\end{tabular}
\end{center}


共 $36$ 种等可能结果。

\textbf{和等于 $7$ 的组合}： $(1,6), (2,5), (3,4), (4,3), (5,2), (6,1)$ ，共 $6$ 个。

\[
P(\text{和} = 7) = \frac{6}{36} = \frac{1}{6}
\]


\end{understand}

\begin{workedexamples}


\textbf{例 4.} 同时掷两枚骰子，求：
（1）点数之和为偶数的概率。
（2）点数之和大于 $10$ 的概率。

\textbf{解：}

利用上面的表格，总共 $36$ 种等可能结果。

（1）点数之和为偶数的组合——两个都是奇数或两个都是偶数：

两奇：骰子 1 取 $\{1,3,5\}$ ，骰子 2 取 $\{1,3,5\}$ ，共 $3 \times 3 = 9$ 种。

两偶：骰子 1 取 $\{2,4,6\}$ ，骰子 2 取 $\{2,4,6\}$ ，共 $3 \times 3 = 9$ 种。

合计 $9 + 9 = 18$ 种。

\[
P(\text{和为偶数}) = \frac{18}{36} = \frac{1}{2}
\]


（2）和大于 $10$ 即和为 $11$ 或 $12$ ：

和 $= 11$ ： $(5,6), (6,5)$ ， $2$ 种。

和 $= 12$ ： $(6,6)$ ， $1$ 种。

合计 $3$ 种。

\[
P(\text{和} > 10) = \frac{3}{36} = \frac{1}{12}
\]


\textbf{例 5.} 袋中有标号为 $1, 2, 3$ 的三个小球。先摸一个记下号码放回，再摸一个。用列表法求两次摸到的号码之和为 $4$ 的概率。

\textbf{解：}

列表：

\begin{center}
\begin{tabular}{llll}
\toprule
第 1 次 \ 第 2 次 & 1 & 2 & 3 \\
\midrule
\textbf{1} & (1,1) & (1,2) & (1,3) \\
\textbf{2} & (2,1) & (2,2) & (2,3) \\
\textbf{3} & (3,1) & (3,2) & (3,3) \\
\bottomrule
\end{tabular}
\end{center}


共 $9$ 种等可能结果。

和为 $4$ 的组合： $(1,3), (2,2), (3,1)$ ，共 $3$ 种。

\[
P(\text{和} = 4) = \frac{3}{9} = \frac{1}{3}
\]


\end{workedexamples}

\begin{keytakeaway}


\begin{itemize}
  \item 列表法适用于有\textbf{两步}实验（或两个因素）的情况。
  \item 列出所有组合后，数出满足条件的个数即可。
  \item 列表法的优势在于\textbf{不重不漏}地列出所有可能结果。
\end{itemize}


\end{keytakeaway}



\section*{用树形图求概率}


\begin{explore}


小明衣柜里有 $2$ 件上衣（白色、蓝色）和 $3$ 条裤子（黑色、灰色、卡其色），他随机各选一件。所有可能的搭配有哪些？选到白色上衣和灰色裤子的概率是多少？

列表法当然可以，但如果涉及\textbf{三步或更多步}操作，表格就画不下了。这时\textbf{树形图}更方便。

\end{explore}

\begin{understand}


\textbf{树形图}从起点出发，每一步的每个选择画一根"树枝"，一直分叉下去，直到所有步骤完成。每条从起点到末端的完整路径就代表一个可能的结果。

小明的穿搭树形图：

```
             ┌── 黑裤 → (白,黑)
     ┌─ 白衣 ┼── 灰裤 → (白,灰)
     │       └── 卡其裤 → (白,卡其)
起点 ─┤
     │       ┌── 黑裤 → (蓝,黑)
     └─ 蓝衣 ┼── 灰裤 → (蓝,灰)
             └── 卡其裤 → (蓝,卡其)
```

共 $2 \times 3 = 6$ 种等可能结果。

\[
P(\text{白衣、灰裤}) = \frac{1}{6}
\]


\end{understand}

\begin{workedexamples}


\textbf{例 6.} 甲、乙、丙三名同学竞选正副班长各一名（一人只能担任一个职位）。用树形图求甲当选正班长的概率。

\textbf{解：}

用树形图列出所有可能的选举结果（正班长，副班长）：

```
         ┌─ 乙 → (甲,乙)
  ┌─ 甲 ─┤
  │      └─ 丙 → (甲,丙)
  │      ┌─ 甲 → (乙,甲)
起点┼─ 乙 ─┤
  │      └─ 丙 → (乙,丙)
  │      ┌─ 甲 → (丙,甲)
  └─ 丙 ─┤
         └─ 乙 → (丙,乙)
```

共 $6$ 种等可能结果。

甲当选正班长的结果： $(甲,乙)$ 和 $(甲,丙)$ ，共 $2$ 种。

\[
P(\text{甲当正班长}) = \frac{2}{6} = \frac{1}{3}
\]


\textbf{例 7.} 甲袋中有 $2$ 个球（红球、白球），乙袋中有 $3$ 个球（红球、白球、黄球）。从每个袋中各随机摸出 $1$ 个球，求摸到的 $2$ 个球颜色相同的概率。

\textbf{解：}

树形图：

```
            ┌─ 红₂ → (红₁,红₂) ✓ 颜色相同
    ┌─ 红₁ ─┼─ 白₂ → (红₁,白₂)
    │       └─ 黄₂ → (红₁,黄₂)
起点 ┤
    │       ┌─ 红₂ → (白₁,红₂)
    └─ 白₁ ─┼─ 白₂ → (白₁,白₂) ✓ 颜色相同
            └─ 黄₂ → (白₁,黄₂)
```

共 $2 \times 3 = 6$ 种等可能结果。

颜色相同的结果： $(红₁,红₂)$ 和 $(白₁,白₂)$ ，共 $2$ 种。

\[
P(\text{颜色相同}) = \frac{2}{6} = \frac{1}{3}
\]


\end{workedexamples}

\begin{keytakeaway}


\begin{itemize}
  \item 树形图适用于\textbf{多步}实验，能清晰展示每一步的所有选择。
  \item 树形图的"路径总数"就是所有等可能结果的总数。
  \item 列表法适合两步实验，树形图适合两步或更多步的实验。
\end{itemize}


\end{keytakeaway}



\section*{实验概率与理论概率}


\begin{explore}


小红做了一个实验：抛一枚图钉 $100$ 次，记录针尖朝上的次数为 $62$ 次。

抛图钉和抛硬币不同——图钉的形状不对称，我们无法像硬币那样直接用"等可能"来算概率。这时该怎么办？

\end{explore}

\begin{understand}


\textbf{理论概率}：通过分析实验的对称性或等可能性，用公式\textbf{计算}出的概率。
\begin{itemize}
  \item 例：抛均匀硬币， $P(\text{正面}) = \frac{1}{2}$ 。
\end{itemize}


\textbf{实验概率}（也叫经验概率）：通过大量实验，用频率来\textbf{估计}的概率。
\begin{itemize}
  \item 例：抛图钉 $100$ 次，针尖朝上 $62$ 次，则估计 $P(\text{针尖朝上}) \approx \frac{62}{100} = 0.62$ 。
\end{itemize}


\textbf{什么时候用实验概率？}
\begin{itemize}
  \item 当各结果\textbf{不等可能}时（如图钉、不规则物体）。
  \item 当问题太复杂，无法理论计算时。
  \item 实验次数越多，实验概率越接近真实概率。
\end{itemize}


\textbf{什么时候用理论概率？}
\begin{itemize}
  \item 当各结果\textbf{等可能}且容易列举时（如骰子、硬币、扑克牌）。
\end{itemize}


\end{understand}

\begin{workedexamples}


\textbf{例 8.} 小刚制作了一个不均匀的转盘，分为红、蓝、绿三个区域。他转了 $200$ 次，记录结果如下：

\begin{center}
\begin{tabular}{llll}
\toprule
颜色 & 红 & 蓝 & 绿 \\
\midrule
次数 & 90 & 70 & 40 \\
\bottomrule
\end{tabular}
\end{center}


估计转到红色的概率约是多少？

\textbf{解：}

\[
P(\text{红色}) \approx \frac{90}{200} = 0.45
\]


因为转盘不均匀，不能直接用面积比来计算理论概率，所以用实验频率 $0.45$ 来估计概率。

\end{workedexamples}

\begin{keytakeaway}


\begin{itemize}
  \item 能用理论方法算的优先用理论概率（精确）；不能算的用实验概率（近似）。
  \item 实验概率的准确度取决于实验次数——次数越多越准。
\end{itemize}


\end{keytakeaway}



\section*{概率的简单应用——游戏公平性}


\begin{explore}


小明和小红玩一个游戏：同时抛两枚硬币，如果两枚都是正面，小明赢；否则小红赢。小红觉得这个游戏"不太公平"，她说得对吗？

\end{explore}

\begin{understand}


\textbf{判断游戏公平性}：如果各玩家\textbf{获胜的概率相等}，则游戏是公平的；否则不公平。

列出同时抛两枚硬币的所有结果（用列表法）：

\begin{center}
\begin{tabular}{lll}
\toprule
硬币 1 \ 硬币 2 & 正 & 反 \\
\midrule
\textbf{正} & (正,正) & (正,反) \\
\textbf{反} & (反,正) & (反,反) \\
\bottomrule
\end{tabular}
\end{center}


共 $4$ 种等可能结果。

\[
P(\text{小明赢}) = P(\text{两枚都正}) = \frac{1}{4}
\]


\[
P(\text{小红赢}) = 1 - \frac{1}{4} = \frac{3}{4}
\]


$\frac{1}{4} \neq \frac{3}{4}$ ，所以这个游戏\textbf{不公平}——小红获胜的概率是小明的 $3$ 倍！小红说得对。

\textbf{如何改成公平游戏？}

方案一：改规则为"两枚相同（都正或都反）小明赢，两枚不同小红赢"。

\[
P(\text{两枚相同}) = \frac{2}{4} = \frac{1}{2}, \quad P(\text{两枚不同}) = \frac{2}{4} = \frac{1}{2}
\]


各 $\frac{1}{2}$ ，公平。

\end{understand}

\begin{workedexamples}


\textbf{例 9.} 甲乙两人玩转盘游戏。转盘被分成 $3$ 个相等的扇形区域，分别标有 $1, 2, 3$ 。规则：转到 $1$ 甲赢，转到 $2$ 或 $3$ 乙赢。这个游戏公平吗？如果不公平，请修改规则使其公平。

\textbf{解：}

\[
P(\text{甲赢}) = \frac{1}{3}, \quad P(\text{乙赢}) = \frac{2}{3}
\]


$\frac{1}{3} \neq \frac{2}{3}$ ，不公平。

修改方案：把转盘改为 $2$ 个相等区域，或者连续转 $2$ 次，规定" $2$ 次都一样甲赢，不一样乙赢"。

转 $2$ 次的方案分析：

\begin{center}
\begin{tabular}{llll}
\toprule
第 1 次 \ 第 2 次 & 1 & 2 & 3 \\
\midrule
\textbf{1} & (1,1) ✓ & (1,2) & (1,3) \\
\textbf{2} & (2,1) & (2,2) ✓ & (2,3) \\
\textbf{3} & (3,1) & (3,2) & (3,3) ✓ \\
\bottomrule
\end{tabular}
\end{center}


$P(\text{两次相同}) = \frac{3}{9} = \frac{1}{3}$ ， $P(\text{两次不同}) = \frac{6}{9} = \frac{2}{3}$ 。仍然不公平。

更简单的修改方案：将转盘改为 $2$ 个相等区域（标 $1$ 和 $2$ ），转到 $1$ 甲赢，转到 $2$ 乙赢。这样 $P(\text{甲赢}) = P(\text{乙赢}) = \frac{1}{2}$ ，公平。

\textbf{例 10.} 盒子里有 $3$ 张卡片，分别写着数字 $1, 2, 3$ 。小华和小李各随机抽一张（先抽不放回），两张卡片上的数字之和为奇数则小华赢，为偶数则小李赢。这个游戏公平吗？

\textbf{解：}

用树形图列出所有结果（不放回抽取）：

```
         ┌─ 2 → (1,2) 和=3 奇 ✓
  ┌─ 1 ─┤
  │      └─ 3 → (1,3) 和=4 偶
  │      ┌─ 1 → (2,1) 和=3 奇 ✓
起点┼─ 2 ─┤
  │      └─ 3 → (2,3) 和=5 奇 ✓
  │      ┌─ 1 → (3,1) 和=4 偶
  └─ 3 ─┤
         └─ 2 → (3,2) 和=5 奇 ✓
```

共 $6$ 种等可能结果。和为奇数的有 $4$ 种，和为偶数的有 $2$ 种。

\[
P(\text{小华赢}) = \frac{4}{6} = \frac{2}{3}, \quad P(\text{小李赢}) = \frac{2}{6} = \frac{1}{3}
\]


不公平，小华获胜的概率是小李的 $2$ 倍。

\end{workedexamples}

\begin{keytakeaway}


\begin{itemize}
  \item 游戏公平的标准是各玩家获胜的概率\textbf{相等}。
  \item 判断公平性的步骤：列出所有等可能结果 $\to$ 计算各方获胜概率 $\to$ 比较。
  \item 概率不等时，可以通过调整规则使游戏变公平。
\end{itemize}


\end{keytakeaway}

\begin{practice}


\textbf{1.} 一个不透明袋子里有 $5$ 个球，其中 $3$ 个红球、 $2$ 个白球。随机摸出 $1$ 个球。
（1）摸到红球的概率是多少？
（2）如果再往袋中加入若干个白球后，使得摸到红球的概率为 $\frac{1}{4}$ ，需要加入几个白球？

\textbf{2.} 同时掷两枚骰子，用列表法求：
（1）两枚点数相同的概率。
（2）点数之和为 $9$ 的概率。
（3）至少有一枚为 $6$ 的概率。

\textbf{3.} 小刚有 $3$ 件不同的上衣和 $2$ 条不同的裤子。
（1）用树形图列出所有可能的搭配。
（2）如果随机选取一套搭配，求选中某一特定搭配的概率。

\textbf{4.} 甲乙两人各掷一枚骰子，点数大的一方获胜，相同则平局。
（1）用列表法求甲获胜的概率。
（2）这个游戏公平吗？

\textbf{5.} 一个袋子里有 $2$ 个红球和 $1$ 个蓝球。不放回地连续摸 $2$ 个球。
（1）用树形图列出所有可能结果。
（2）求两个球颜色相同的概率。
（3）求至少摸到一个红球的概率。

\end{practice}



\begin{answer}


\textbf{1.}

（1） $P(\text{红球}) = \frac{3}{5}$ 。

（2）设加入 $n$ 个白球，则总球数 $= 5 + n$ ，红球仍为 $3$ 个。

\[
\frac{3}{5 + n} = \frac{1}{4}
\]


$3 \times 4 = 5 + n$ ， $12 = 5 + n$ ， $n = 7$ 。需要加入 $7$ 个白球。

\textbf{2.}

共 $36$ 种等可能结果。

（1）两枚点数相同： $(1,1), (2,2), (3,3), (4,4), (5,5), (6,6)$ ，共 $6$ 种。

\[
P(\text{相同}) = \frac{6}{36} = \frac{1}{6}
\]


（2）和为 $9$ ： $(3,6), (4,5), (5,4), (6,3)$ ，共 $4$ 种。

\[
P(\text{和}=9) = \frac{4}{36} = \frac{1}{9}
\]


（3）至少一枚为 $6$ ：用对立事件法更简便。

"两枚都不是 $6$ "：每枚有 $5$ 种可能（ $1 \sim 5$ ），共 $5 \times 5 = 25$ 种。

\[
P(\text{至少一枚为 } 6) = 1 - \frac{25}{36} = \frac{11}{36}
\]


\textbf{3.}

（1）设上衣为 $A, B, C$ ，裤子为 $P, Q$ 。

```
         ┌─ P → (A,P)
  ┌─ A ─┤
  │      └─ Q → (A,Q)
  │      ┌─ P → (B,P)
起点┼─ B ─┤
  │      └─ Q → (B,Q)
  │      ┌─ P → (C,P)
  └─ C ─┤
         └─ Q → (C,Q)
```

共 $3 \times 2 = 6$ 种搭配。

（2） $P(\text{某一特定搭配}) = \frac{1}{6}$ 。

\textbf{4.}

（1）从 $36$ 种结果中，甲点数 $>$ 乙点数的情况：

\begin{itemize}
  \item 甲 $2$ ，乙 $1$ ： $1$ 种
  \item 甲 $3$ ，乙 $1$ 或 $2$ ： $2$ 种
  \item 甲 $4$ ，乙 $1, 2, 3$ ： $3$ 种
  \item 甲 $5$ ，乙 $1, 2, 3, 4$ ： $4$ 种
  \item 甲 $6$ ，乙 $1, 2, 3, 4, 5$ ： $5$ 种
\end{itemize}


合计 $1 + 2 + 3 + 4 + 5 = 15$ 种。

\[
P(\text{甲赢}) = \frac{15}{36} = \frac{5}{12}
\]


（2）由对称性， $P(\text{乙赢}) = \frac{15}{36} = \frac{5}{12}$ 。

$P(\text{平局}) = \frac{6}{36} = \frac{1}{6}$ 。

验证： $\frac{5}{12} + \frac{5}{12} + \frac{1}{6} = \frac{5}{12} + \frac{5}{12} + \frac{2}{12} = 1$ ✓

$P(\text{甲赢}) = P(\text{乙赢}) = \frac{5}{12}$ ，这个游戏是\textbf{公平的}。

\textbf{5.}

（1）设 $2$ 个红球为红 $_1$ 、红 $_2$ ，蓝球为蓝。

```
          ┌─ 红₂ → (红₁,红₂)
  ┌─ 红₁ ─┤
  │       └─ 蓝 → (红₁,蓝)
  │       ┌─ 红₁ → (红₂,红₁)
起点┼─ 红₂ ─┤
  │       └─ 蓝 → (红₂,蓝)
  │       ┌─ 红₁ → (蓝,红₁)
  └─ 蓝 ──┤
          └─ 红₂ → (蓝,红₂)
```

共 $6$ 种等可能结果。

（2）两球颜色相同的只有两个都是红球： $(红_1,红_2)$ 和 $(红_2,红_1)$ ，共 $2$ 种。

\[
P(\text{颜色相同}) = \frac{2}{6} = \frac{1}{3}
\]


（3）至少摸到一个红球——用对立事件。"没有红球"意味着两个都是蓝球，但只有 $1$ 个蓝球，不可能连续摸出 $2$ 个蓝球。

\[
P(\text{至少一个红球}) = 1 - 0 = 1
\]


即必然至少摸到一个红球。

\end{answer}
